\section{Die Frage}

$\theta = 2/9$ ist empirisch auf sieben Stellen bestätigt (A110). Dieses
Dokument zeigt dreierlei: \emph{wo} der Wert sitzt (allein im dynamischen
Phasenkanal, nicht im Betrag), \emph{dass} er sich exakt und parameterfrei aus
der ikosaedrischen Verdopplungsgeometrie ergibt, und \emph{was} als
Ableitungsschritt noch offen ist (die Brücke zu den Massenverhältnissen).

Der zentrale Befund ist positiv: zwei voneinander unabhängige Wege -- der
phänomenologische Koide-Winkel aus den gemessenen Massen und die rein
geometrische Umverteilung der Moden -- treffen exakt dieselbe rationale Zahl
$2/9$, ohne voneinander zu wissen.

\section{Der Betragskanal ist blind}

\textbf{[K]} Mit den Amplituden $a_k(\theta) = 1 + \sqrt2\cos(\theta+2\pi k/3)$
ist der Koide-Invariant
\[
  Q(\theta) = \frac{\sum_k a_k^2}{\bigl(\sum_k a_k\bigr)^2} = \frac{6}{9} = \frac23
  \qquad\text{für \emph{jedes} } \theta .
\]

Der Beweis ist zwei Zeilen lang: drei um $120^\circ$ versetzte Kosinus summieren
zu null, also $\sum_k a_k = 3$; und mit $\sum_k\cos^2 = 3/2$ folgt
$\sum_k a_k^2 = 6$.

\textbf{[K]} Daraus folgt eine Aussage, die man festhalten muss, weil sie die
Reichweite von A110 begrenzt: \emph{die Koide-Relation sagt über $\theta$
nichts}. $Q=2/3$ gilt für jeden Winkel. Der Wert $2/9 = 0{,}2222$ liegt
zudem \emph{unterhalb} des Wertebereichs $[1/3,1]$, den $Q$ bei positiven
Massen durchläuft -- $2/9$ ist also weder Grenzwert noch Extremum noch
Fixpunkt des symmetrischen Funktionals.

Im Betragskanal ist $2/9$ vollständig unsichtbar.

\section{Die Eliminationskette}

\textbf{[K]} Systematisch durchgeprüft, wo $\theta$ \emph{nicht} lebt. Die
Orbit-Glieder $\{1/9,\,2/9,\,4/9\}$ sind zu unterscheiden; die Frage ist,
welche Größe das leistet.

\begin{center}
\begin{tabular}{lll}
\hline
Kanal & Verhalten & unterscheidet? \\
\hline
Betrag $|B_\theta|$ & $=1$ bis $10^{-15}$ für alle $\theta$ & nein \\
Windungszahl & $=3$ exakt für alle $\theta$ & nein \\
symmetrisches Funktional & identisch für alle $\theta$ & nein \\
stetiges Phasenprofil & $0{,}093$ / $0{,}171$ / $0{,}249$ & \textbf{ja} \\
\hline
\end{tabular}
\end{center}

\textbf{[K]} Bemerkenswert ist die zweite Zeile. Die ganzzahlige Windung -- also
gerade die topologische Größe, die nach A030 den Informationsbegriff der Theorie
trägt -- \emph{trennt $2/9$ nicht} von $1/9$ und $4/9$. Der Grund ist,
dass $2/(9\pi)$ irrational ist; $\theta$ ist kein flacher topologischer
Invariant.

\textbf{[K]} Übrig bleibt genau ein Kanal: das stetige, dynamische Phasenprofil.
Dort und nur dort ist $2/9$ von seinen Orbit-Geschwistern unterscheidbar.

\section{Erzwungen oder kontingent}

\textbf{[K]} Damit stellt sich die entscheidende Frage: ist $2/9$ erzwungen
wie die Lichtgeschwindigkeit, oder kontingent?

Der Test ist das Variationsverhalten. Koppelt man ein zweites Triplett mit
Stärke $\delta$ an, wandert das Extremum vom $Z_3$-Punkt weg:
$\theta^*$ steigt glatt und monoton von $0{,}099$ bei $\delta = 0{,}10$
über $0{,}223$ bei $\delta = 0{,}24$ bis $0{,}501$ bei $\delta = 1{,}0$.
Der erreichte Bereich $[0{,}05;\,0{,}52]$ enthält $2/9$ \emph{im Inneren und
unausgezeichnet}.

\textbf{[K]} Dieser Test setzt $\delta$ als \emph{freien} Kopplungsparameter an.
Unter dieser Annahme läge $2/9$ unausgezeichnet im erreichten Bereich, und der
Wert wäre kontingent. Die Annahme ist aber nicht die der Theorie: die
ikosaedrische Geometrie legt $\delta$ fest (nächster Abschnitt). Sobald $\delta$
nicht mehr frei ist, sondern der totale Leak $(7-3\varphi)/9$ derselben
Umverteilung, die auch $\theta$ trägt, entfällt die Kontingenz. Der
Variationstest zeigt also, dass $\theta$ \emph{im dynamischen Kanal} lebt -- die
Frage erzwungen/kontingent entscheidet er nicht, weil er $\delta$ als frei
unterstellt, was die Geometrie nicht tut.

\section{Die geometrische Quelle: $C_3$ in $A_5$}

\textbf{[B]} Der dynamische Kanal hat eine geometrische Quelle. Die drei
Lepton-Moden sind die Fourier-Eigenmoden des $Z_3$-Zirkulanten: die triviale
Mode $v_0 = \tfrac{1}{\sqrt3}(1,1,1)$ und zwei nicht-triviale, galois-konjugierte
Moden. Die Verdopplung stellt die $C_3$-Struktur in die ikosaedrische Gruppe
$A_5$; deren fünfzählige Erzeugende ist eine Drehung $R_5$ um $72^\circ$ um eine
ikosaedrische Fünffach-Achse. Weil die drei- und die fünfzählige Achse
verschieden orientiert sind, mischt $R_5$ die $C_3$-Moden.

\textbf{[B]} Die volle Umverteilung der Elektron-Mode hat die Gewichte
$p_j = |\langle v_j\,|\,R_5\,v_e\rangle|^2$:
\[
  \underbrace{p_0 = \tfrac29}_{\text{trivial}},\qquad
  p_2 = \tfrac{5-3\varphi}{9},\qquad
  p_1 = \tfrac{2+3\varphi}{9},\qquad
  p_0+p_1+p_2 = 1,
\]
mit $\varphi = \tfrac{1+\sqrt5}{2}$. \emph{Alle drei Gewichte tragen den Nenner
$9 = 3^2$} -- dieselbe $Z_3$-Struktur, in der $2/9$ als Orbit-Adresse sitzt. Der
Anteil, der unter der Fünffach-Drehung in die triviale Mode übergeht, ist
\[
  p_0 = \bigl|\langle v_0\,|\,R_5\,v_e\rangle\bigr|^2 = \frac29
\]
exakt und parameterfrei (numerisch auf 40 Stellen bestätigt). In diesen Wert
geht kein Massenverhältnis, keine Koide-Parametrisierung und kein $2/9$ ein.

\textbf{[K]} Der Wert ist ikosaeder-\emph{spezifisch}, nicht generisch:
\begin{itemize}
  \item $200$ zufällige Drehachsen bei $72^\circ$ geben den Leak im Bereich
    $[0{,}036;\,0{,}452]$; \emph{keine einzige} trifft $2/9$. Nur die
    ikosaedrischen Fünffach-Achsen liefern ihn.
  \item Der Winkel-Scan zeigt $2/9$ genau bei $72^\circ$; die $n$-zählige Reihe
    liefert $2/9$ ausschließlich bei $n = 5$ (bei $n=3$ steht $2/5$, bei
    $144^\circ$ steht $4/9$).
  \item Von den sechs Fünffach-Achsen geben drei den Wert $2/9$, drei den Wert
    $4/9$; die Einbettung wählt die nächstliegende Achse und damit $2/9$. Das ist
    eine diskrete Zweifachwahl, kein freier Parameter.
\end{itemize}

\textbf{[B]} Dieselbe Umverteilung liefert auch die zweite Größe: der totale
Leak $p_0 + p_2 = (7-3\varphi)/9 = 0{,}2384$ ist der parameterfreie Wert der
Kopplung $\delta$ und stimmt mit dem aus $\theta$ rückgerechneten
$\delta^* = 0{,}2389$ auf $0{,}18\,\%$ überein. Winkel \emph{und} Amplitude
entstammen also einer einzigen Fünffach-Umverteilung; die frühere getrennte
Behandlung als zwei Harmonische war die zu enge Sicht.

\textbf{[K]} Damit fallen zwei unabhängige Wege exakt zusammen: der
phänomenologische Koide-Winkel (aus den Massen, ohne Geometrie) und die
geometrische Umverteilung (aus $A_5$, ohne Massen und ohne $2/9$-Input) nennen
beide $2/9$. Das ist ein Konvergenzbeleg. $2/9$ ist ein Invariant der
Übersetzung zwischen der Torusgeometrie und der Spektraldarstellung -- er gehört
der gemeinsamen Struktur, nicht einer Darstellung, ähnlich wie $\hbar$ und $c$
Umrechnungsfaktoren sind und keine ontologischen Primitive.

\textbf{[B]} Der Anschluss an den dynamischen Kanal schließt sich hier: die
Fünffach-Drehung $R_5$ setzt eine Phase $\pi/2$ -- genau den Wert $\chi = \pi/2$,
der oben den einen unterscheidenden Kanal markiert. Die geometrische Quelle und
der dynamische Ort sind dasselbe.

\section{Prüfskript}

\begin{itemize}
  \item $\theta$-Unabhängigkeit von $Q$ (exakt, für beliebige Winkel), die
    Eliminationskette über die drei blinden Kanäle, und die Kontingenzprobe
    $\theta^*(\delta)$ mit dem Nachweis, dass $2/9$ im erreichten Bereich
    unausgezeichnet liegt:
    \texttt{python/A\_Serie\_Skripte/a120\_theta.py}
\end{itemize}


\section{Zeichenerklärung}

Symbole die in der gesamten A-Serie verwendet werden, sind in A010 erklärt.
Spezifisch für dieses Dokument:

\begin{center}
\renewcommand{\arraystretch}{1.3}
{\small\begin{tabular}{lp{9cm}}
\toprule
Symbol & Bedeutung \\
\midrule
$\theta = 2/9$ & Geometrischer Fixpunkt-Winkel; dynamischer Locus auf $T^4$ \\
$p_0 = 2/9$ & Exakter $\theta$-Wert als $C_3$-in-$A_5$-Identität \\
$\Omega_\text{res}$ & Resonanzfrequenz der $\theta$-Mode \\
$\Delta\omega$ & Frequenzabstand zur nächsten Mode \\
\bottomrule
\end{tabular}}
\renewcommand{\arraystretch}{1.0}
\end{center}

\section{Was offen bleibt}

\textbf{Erstens ist der Wert $2/9$ hergeleitet, nicht offen.} Er folgt exakt und
parameterfrei aus der ikosaedrischen Umverteilung und fällt unabhängig mit dem
phänomenologischen Koide-Winkel zusammen. Die frühere Einordnung als
strukturell unbegründet ist damit überholt.

\textbf{Zweitens ist $\delta$ nicht frei.} Es ist der totale Leak
$(7-3\varphi)/9$ derselben Umverteilung, die $\theta$ trägt. Der
Kontingenz-Eindruck entstand allein daraus, $\delta$ als freien Parameter zu
behandeln; die Geometrie legt ihn fest.

\textbf{Drittens bleibt eine Ableitungskante offen -- und sie ist genau benannt.}
Der explizite Mechanismus, der die geometrischen Umverteilungsgewichte
$(2/9,\,(5-3\varphi)/9,\,(2+3\varphi)/9)$ mit den physikalischen
Massenverhältnissen verbindet, ist nicht ausgeschrieben. Diese Brücke muss nicht
über Koide laufen; sie kann direkt über die Massen oder die Modenstruktur
geführt werden. Das ist der Gegenstand der Weiterarbeit -- und ausdrücklich kein
Einwand gegen den Konvergenzbeleg, sondern die nächste Stufe.

\textbf{Viertens ist die zugrunde liegende Verdopplungsstruktur $A_5$
vorausgesetzt.} Der Wert $2/9$ ist innerhalb der etablierten ikosaedrischen
Verdopplung exakt; ob diese Verdopplung selbst aus den drei Setzungen von A020
zwingend folgt oder als Struktur hinzutritt, ist die verbleibende strukturelle
Frage -- eine über die Zähligkeit hinausgehende, die den Wert $2/9$ aber nicht
mehr berührt.

\section{Ersetzt}

Dieses Dokument nimmt auf: Dok.~291 (dynamischer Ort), Dok.~293 (ikosaedrischer
Zugang), Dok.~286 (dynamischer Sektor), Dok.~290 (Betrag/Phase, soweit $\theta$ betreffend). Eingearbeitet: P42, P35.

\section{Verwendung von KI-Werkzeugen}

Dieses Dokument ist unter Verwendung KI-gestützter Werkzeuge entstanden. Die
Fragestellungen, die Setzungen und alle inhaltlichen Entscheidungen stammen vom
Autor; die Werkzeuge formulieren aus, rechnen nach und erstellen Prüfskripte.
Die Verantwortung für Inhalt und Fehler liegt vollständig beim Autor. Der
ausführliche Wortlaut steht in A010.
